\section{Lecture 6, Part 1: Model Families \& Hybrid Retrieval}

\begin{frame}
\vfill
\centering
{\color{RLBlue}\large Lecture 6}

\vspace{0.3em}
{\color{RLBlue}\Large\bfseries Part 1}

\vspace{0.6em}
{\Large Model Families \& Hybrid Retrieval}
\vfill
\end{frame}

\begin{frame}{Learning Goals}
\begin{itemize}
  \item Read an MTEB leaderboard and say what it does \emph{not} tell you about your corpus (Muennighoff et al., 2022).
  \pause
  \item Explain what late interaction buys over a single vector, and what it costs to store (Khattab \& Zaharia, 2020).
  \pause
  \item Combine a sparse and a dense ranking, and justify why hybrid is the default rather than a compromise.
\end{itemize}
\end{frame}

% ─── Model families ─────────────────────────────────────────────────────────
\begin{frame}{Why Are There So Many Embedding Models?}
\begin{itemize}
  \item \textbf{Why:} the bi-encoder recipe from Part 0 -- encode separately, train contrastively -- is not one model. It is a template, and every family fills it differently.
  \pause
  \item \textbf{E5} (Wang et al., 2022) trains on weakly-supervised pairs mined at scale, and asks you to prefix inputs with \texttt{query:} or \texttt{passage:} so one model can play both towers.
  \pause
  \item \textbf{GTE} (Li et al., 2023) uses multi-stage contrastive training, broad pretraining followed by supervised fine-tuning.
  \pause
  \item \textbf{BGE-M3} (Chen et al., 2024) is multilingual and \emph{multi-functional}: one model emitting a dense vector, sparse term weights and multi-vector output at once -- so hybrid retrieval needs one model rather than three.
  \pause
  \item \textbf{Hosted APIs} (OpenAI's embedding models among them) trade control for convenience: no weights to run, but no offline mode and no ability to inspect what changed between versions.
\end{itemize}
\end{frame}

\begin{frame}{How to Read MTEB}
\begin{rldefinition}
\textbf{MTEB} -- the Massive Text Embedding Benchmark (Muennighoff et al., 2022): one leaderboard covering eight task families, including retrieval, clustering, classification and semantic similarity, across dozens of datasets and many languages. \emph{Why it exists:} a model tuned for semantic similarity can be mediocre at retrieval, and a single average hides that.
\end{rldefinition}
\pause
\begin{itemize}
  \item \textbf{Read the retrieval column, not the average.} You are building a retriever. The other seven task families are someone else's problem.
  \pause
  \item \textbf{Check the sequence limit.} A model truncating at 512 tokens will silently cut long articles in half, and the half it keeps may not be the half that answers the question.
  \pause
  \item \textbf{Check the cost of a point.} A model one point better on a leaderboard but four times the dimension costs four times the index memory -- Lecture 7 puts a number on that.
\end{itemize}
\end{frame}

\begin{frame}{Why MTEB Cannot Answer Your Question}
\begin{itemize}
  \item \textbf{Why:} a leaderboard measures the datasets it contains. Yours is not one of them.
  \pause
  \item \textbf{The domain gap.} Legal drafting is unlike web text in every way that matters for retrieval: fixed phrasing, heavy cross-reference, long sentences, and words that are ordinary elsewhere and technical here.
  \pause
  \item \textbf{The language gap.} BGE-M3 reports support for over a hundred languages. Whether that support is \emph{good} for Uzbek is not something any published evaluation currently answers -- no Uzbek-specific retrieval benchmark exists.
  \pause
  \item \textbf{So measure it.} A team that reports recall@$k$ for two or three embedding models on an Uzbek legal corpus has produced a number nobody has published. That is a genuine contribution, and it is a sprint's work.
\end{itemize}
\end{frame}

% ─── Late interaction ───────────────────────────────────────────────────────
\begin{frame}{Why Does One Vector Per Article Lose Information?}
\begin{itemize}
  \item \textbf{Why:} a bi-encoder compresses a whole article into $d$ numbers \emph{before} it knows the question. Everything the question cared about had to survive that compression.
  \pause
  \item An article covering working hours, the employer's scheduling duty, and exceptions becomes one point. A question about any one of those three is compared against the average of all three.
  \pause
  \item A cross-encoder has no such problem -- it reads the question first -- but Part 0 showed it cannot be a first stage.
  \pause
  \item \textbf{Late interaction} is the compromise: keep one vector \emph{per token}, still encoded independently, and postpone the comparison to a cheap operation over those vectors.
\end{itemize}
\end{frame}

\begin{frame}{ColBERT: Late Interaction}
\begin{rldefinition}
\textbf{Late interaction} (Khattab \& Zaharia, 2020) -- encode query and document separately into \emph{per-token} vectors, then score with \textbf{MaxSim}: for each query token take its best match anywhere in the document, and add those up:
\vspace{-0.3em}
\[
  \mathrm{score}(Q,D) \;=\; \sum_{i \in Q} \max_{j \in D} \; \mathbf{q}_i \cdot \mathbf{d}_j
\]
\vspace{-0.6em}

\emph{Here:} the query token ``hours'' finds ``hours'' in Article 182 directly, instead of hoping the article's one average vector leaned that way.
\end{rldefinition}
\pause
\begin{itemize}
  \item Still a bi-encoder where it counts: the document side never sees the query, so it is precomputed once. Only the \emph{comparison} is deferred -- hence ``late''.
\end{itemize}
\end{frame}

\begin{frame}{What Late Interaction Costs}
\begin{itemize}
  \item \textbf{The arithmetic:} the capstone corpus is $7{,}368$ articles and about $971{,}000$ words -- roughly $132$ words per article.
  \pause
  \item One vector per \emph{article}, at $d = 1024$ in \texttt{float16}: $7{,}368 \times 1024 \times 2$ bytes $\approx$ \textbf{15\,MB}. That fits anywhere.
  \pause
  \item One vector per \emph{token}, on the same corpus and the same $d$: roughly $1.4$ million vectors $\approx$ \textbf{2.8\,GB} -- about $185\times$ more, for the same text.
  \pause
  \item That is the trade in one line: late interaction buys back the detail a single vector threw away, and pays for it in index size. ColBERTv2 (Santhanam et al., 2022) exists mainly to make that bill smaller through compression.
  \pause
  \item \textbf{For a corpus this size, either is affordable.} Make the choice on measured recall, not on instinct -- and note that $185\times$ becomes decisive the moment the corpus does.
\end{itemize}
\end{frame}

% ─── Hybrid ─────────────────────────────────────────────────────────────────
\begin{frame}{Why Do We Need Both Sparse and Dense?}
\begin{itemize}
  \item \textbf{Why:} the two methods fail at different things, and neither failure is rare.
  \pause
  \item \textbf{Sparse fails on paraphrase.} Part 0's second query scored zero on every article, because the user's words were not the drafter's words.
  \pause
  \item \textbf{Dense fails on exact tokens.} ``Article 182'', a date, a statute number, a rare compound -- these are precisely what a $1024$-dimensional average is worst at preserving, and precisely what a legal query is most likely to contain.
  \pause
  \item \textbf{They also fail independently}, which is the useful part. A query that defeats one is usually handled by the other, so the union of their top-$k$ has markedly better recall than either alone.
\end{itemize}
\end{frame}

\begin{frame}{Combining Two Rankings}
\begin{itemize}
  \item \textbf{The obvious approach does not work.} BM25 returns an unbounded score built from IDF sums; cosine returns something in $[-1,1]$. Adding them means letting the corpus statistics decide the weighting, differently for every query.
  \pause
  \item \textbf{Fuse ranks, not scores.} Reciprocal rank fusion combines lists using only the position a document reached in each -- no calibration, no tuning, nothing to go stale. Lecture 5 covers it as a reranking stage.
  \pause
  \item \textbf{Or score, then rerank.} Take the union of both top-$k$ lists and let a cross-encoder order the survivors. Expensive per document, but there are now only a few dozen.
  \pause
  \item \textbf{BGE-M3's argument} for doing all this in one model is that its dense and sparse outputs come from the same encoder, so they are at least trained together rather than bolted together.
\end{itemize}
\end{frame}

\begin{frame}{Why This Matters: BM25 Is the Number to Beat}
\begin{itemize}
  \item Sprint 2 of the capstone is scored as recall@$k$ against your Sprint 1 baseline. That baseline is a sparse index, and on statute text it will be \emph{good}.
  \pause
  \item \textbf{Expect dense retrieval to lose at first.} An off-the-shelf multilingual model, on formulaic legal text in a language it was barely trained on, against a keyword index on a corpus where the drafter chose the words deliberately -- BM25 winning is the normal result, not a bug in your code.
  \pause
  \item \textbf{The finding is the deliverable.} ``Dense retrieval did not beat BM25 here, and here is the recall table'' is a better report than an unmeasured claim that it did.
  \pause
  \item \textbf{And hybrid is usually the honest answer} -- not a hedge, but the design the two failure modes actually imply.
\end{itemize}
\end{frame}

\begin{frame}{Choosing a Model for the Capstone}
\begin{itemize}
  \item \textbf{Embeddings are a separate decision from generation.} The agent's language model is a hosted API; the embedding model need not be.
  \pause
  \item \textbf{Local is viable here.} Encoding $7{,}368$ articles -- about $1.4$ million tokens -- is a one-time job of roughly an hour on a laptop CPU. After that, queries are free, offline and unversioned by anyone else.
  \pause
  \item \textbf{A hosted embedding API} costs cents for the same corpus, and removes the download. It also means your index silently depends on someone else's model version.
  \pause
  \item \textbf{Whatever you pick, pin it and record it.} An embedding model change invalidates every vector in your index -- and a recall number computed under a different model is not comparable to last sprint's.
\end{itemize}
\end{frame}

\begin{frame}{Summary}
\begin{itemize}
  \item \textbf{BM25} scores rarity times damped frequency, normalised by length. On statute text it is a strong baseline, not a warm-up.
  \pause
  \item Its one fatal gap is \textbf{paraphrase}: a query sharing no content word scores zero everywhere, and zero is unrecoverable.
  \pause
  \item A \textbf{bi-encoder} closes that gap by encoding both sides separately -- the only design where the corpus can be encoded once instead of once per query.
  \pause
  \item Embedding models are \textbf{trained contrastively}: in-batch negatives make the easy comparisons free, hard negatives supply the ones that still teach something.
  \pause
  \item \textbf{Late interaction} recovers per-token detail at roughly $185\times$ the index size on this corpus.
  \pause
  \item \textbf{Hybrid retrieval is the default}, because sparse and dense fail on different queries -- and a leaderboard cannot tell you which failure dominates on \emph{your} corpus.
\end{itemize}
\end{frame}

\begin{frame}{Reading}
\begin{itemize}
  \item Robertson \& Zaragoza (2009). \textit{The Probabilistic Relevance Framework: BM25 and Beyond.} --- the baseline, and where $k_1$ and $b$ come from.
  \pause
  \item Karpukhin et al.\ (2020). \textit{Dense Passage Retrieval for Open-Domain Question Answering.} EMNLP. --- the bi-encoder recipe and hard negatives.
  \pause
  \item Khattab \& Zaharia (2020). \textit{ColBERT: Efficient and Effective Passage Search via Contextualized Late Interaction over BERT.} SIGIR.
  \pause
  \item Chen et al.\ (2024). \textit{BGE M3-Embedding: Multi-Lingual, Multi-Functionality, Multi-Granularity Text Embeddings.} --- the model this project is most likely to use.
  \pause
  \item Muennighoff et al.\ (2022). \textit{MTEB: Massive Text Embedding Benchmark.} --- and its limits on a corpus like yours.
  \pause
  \item \textbf{Next lecture:} those vectors are useless until you can search a million of them in milliseconds -- HNSW, IVF and product quantization.
\end{itemize}
\end{frame}
